\documentclass{article}
\usepackage[utf8]{inputenc}
\usepackage[ngerman]{babel} % deutsch
\usepackage{amsmath}
\usepackage{amsfonts}
\usepackage[colorlinks=true, urlcolor=blue, linkcolor=blue]{hyperref}
\usepackage{tasks}

\title{Mögliche Lösungen zum Reserve Mathe 5 BAC 2017 B-Teil}
\author{Milix}
\date{\today}

\begin{document}

\maketitle

\section*{Aufgabe 1}
Verwenden Sie den Rechner bei den Aufgaben e) und f). \\
Gegeben ist eine Funktion \(f\) durch
\[f(x)=\frac{x^3-12x^2-52x-64}{20x}\]
\begin{tasks}
    \task Zeigen Sie, dass das Schaubild von \(f\) genau eine Asymptote besitzt
    und geben Sie eine Gleichung dieser Asymptote an.
    \task Untersuchen Sie die Bereiche des Steigens und Fallens der Funktion \(f\).
    \task Bestimmen Sie die Koordinaten des Extrempunkts von \(f\).
    \task Zeigen Sie, dass das Schaubild von \(f\) einen Wendepunkt besitzt und
    bestimmen Sie eine Gleichung der Tangente des Schaubilds von \(f\) in
    diesem Punkt.
    \task Im ersten Quadranten befindet sich eine Fläche, die von der x-Achse,
    der Gerade mit der Gleichung \(y=\dfrac{5}{4}x\) und dem Schaubild von \(f\)
    begrenzt wird. \\
    Berechnen Sie den Inhalt dieser Fläche.
    \task Gegeben ist der Punkt \(P(0|-3)\) . Es gibt durch diesen Punkt eine
    Gerade, die Tangente am Schaubild von \(f\) ist. \\
    Ermitteln Sie eine Gleichung für diese Tangente und bestimmen Sie die
    Koordinaten ihres Berührpunkts am Schaubild von \(f\).
\end{tasks}

\subsection*{Lösung}

\newpage
\section*{Aufgabe 2}
Verwenden Sie den Rechner bei der Aufgabe b). \\
In einem 3-dimensionsalen Raum sind gegeben : \\
die Punkte
\[A(1|2|2),\: B(-5|2|10),\: C(1|-1|0),\: D(3|2|4) \]
die Gerade
\[
    d_1 : \begin{pmatrix} x \\ y \\ z \end{pmatrix}
    = \begin{pmatrix} 1 \\ -1 \\ 0 \end{pmatrix}
    + \lambda \begin{pmatrix} 2 \\ 1 \\ -2 \end{pmatrix}, \quad \lambda \in \mathbb{R},
\]
und
\begin{center}
    die Gerade \(d_2\), die durch die Punkte \(A\) und \(B\) geht.
\end{center}
\begin{tasks}
    \task Zeigen Sie, dass sich die Geraden \(d_1\) und \(d_2\) im Punkt \(P(7|2|-6)\) schneiden.

    \task Bestimmen Sie in Grad die Größe des spitzen Winkels, unter dem sich die Geraden \(d_1\) und \(d_2\) schneiden.

    \task Ermitteln Sie eine Koordinatengleichung für die Ebene, auf der die Geraden \(d_1\) und \(d_2\) liegen.

    \task Berechnen Sie den Abstand zwischen dem Punkt \(D\) und der Ebene \(\pi\), die das Dreieck \(APC\) enthält.

    \task Es gibt eine Ebene, die senkrecht auf der Ebene \(\pi\) steht und die die Punkte \(A\) und \(D\) enthält.

    Bestimmen Sie eine Koordinatengleichung dieser Ebene.

    \task Ermitteln Sie eine Gleichung für die Kugel, die die Strecke \([DP]\) als Durchmesser besitzt.
\end{tasks}
\subsection*{Lösung}

\newpage
\section*{Aufgabe 3}
Verwenden Sie den Rechner bei den Aufgaben d), e) und f). \\
Das Unternehmen Fructidoux füllt seine Konfitüre in kleine Gläser.
Fructidoux hat das Ziel, die Bezeichnung „Konfitüre leicht`` zu erhalten. Um
diese Bezeichnung zu erhalten, muss nach den gesetzlichen Vorschriften
der Zuckergehalt pro Gramm Konfitüre zwischen \(0{,}16\) Gramm und
\(0{,}18\) Gramm liegen. \\
Das Unternehmen besitzt die beiden Abfüllanlagen \(F_1\) und \(F_2\).

\begin{itemize}
    \item \(70\,\%\) der Gläser stammen von \(F_1\),
    \item \(30\,\%\) der Gläser stammen von \(F_2\).
    \item \(5\,\%\) der Gläser von \(F_1\) entsprechen nicht den gesetzlichen Vorschriften,
    \item \(1\,\%\) der Gläser von \(F_2\) entsprechen nicht den gesetzlichen Vorschriften.
\end{itemize}

\noindent Der Gesamtproduktion wird zufällig ein Glas entnommen.

\begin{tasks}(1)
    \task Berechnen Sie die Wahrscheinlichkeit, dass dieses Glas den
    gesetzlichen Vorschriften entspricht und von der Abfüllanlage \(F_1\)
    stammt.

    \task Zeigen Sie, dass die Wahrscheinlichkeit, dass dieses Glas den
    gesetzlichen Vorschriften entspricht, den Wert \(0{,}962\) besitzt.

    \task Wenn gegeben ist, dass dieses Glas den gesetzlichen Vorschriften
    entspricht, berechnen Sie die Wahrscheinlichkeit, dass es von der
    Abfüllanlage \(F_2\) stammt.
\end{tasks}

\noindent Der Gesamtproduktion werden zufällig \(1000\) Gläser entnommen.

\begin{tasks}(1)
    \task[d)] Verwenden Sie eine geeignete Wahrscheinlichkeitsverteilung und
    berechnen Sie die Wahrscheinlichkeit, dass mindestens \(950\) dieser
    Gläser den gesetzlichen Vorschriften entsprechen.
\end{tasks}

\noindent Ein Hotel kauft \(50\) Gläser von der Gesamtproduktion.

\begin{tasks}(1)
    \task[e)] Berechnen Sie mit Hilfe einer Poissonverteilung die Wahrscheinlichkeit,
    dass mindestens \(5\) Gläser nicht den gesetzlichen Vorschriften
    entsprechen.
\end{tasks}

\noindent Das Unternehmen erwirbt eine dritte Abfüllanlage \(F_3\). \\
Bei dieser Abfüllanlage folgt der Zuckergehalt pro Gramm der Konfitüre einer
Normalverteilung mit dem Mittelwert \(\mu = 0{,}17\) Gramm und der
Standardabweichung \(\sigma\).
\vspace{0.2cm}
\noindent Man entnimmt zufällig ein Glas aus der Produktion von \(F_3\).
Die Wahrscheinlichkeit, dass dieses Glas den gesetzlichen Vorschriften
entspricht, beträgt \(0{,}90\).

\begin{tasks}(1)
    \task[f)] Bestimmen Sie \(\sigma\).
\end{tasks}

\subsection*{Lösung}

\newpage
\section*{Aufgabe 4}

Eine Folge \((u_n)\) ist gegeben durch:
\[
    \begin{cases}
        u_1 = 0{,}1 \\
        u_{n+1} = 0{,}64\, u_n + 0{,}18 \quad \text{mit } n \in \{1, 2, 3, \dots\}
    \end{cases}
\]

\begin{tasks}(1)
    \task Stellen Sie die Folge \((u_n)\) graphisch dar mit Hilfe einer hinreichenden Anzahl von Punkten \((n \,/\, u_n)\) und leiten Sie daraus eine Vermutung für den Grenzwert der Folge \((u_n)\) her.
    \task Setzen Sie die Konvergenz der Folge \((u_n)\) voraus und bestimmen Sie ihren Grenzwert.
\end{tasks}
\subsection*{Lösung}

\newpage
\section*{Aufgabe 5}
Gegeben sind die komplexen Zahlen \(z_1 = e^{i\frac{\pi}{4}}\) und \(z_2 = 1 + e^{i\frac{\pi}{4}}\).

\begin{tasks}(1)
    \task Stellen Sie \(z_1\) und \(z_2\) in der Gaußschen Zahlenebene dar und bestimmen Sie geometrisch daraus das Argument von \(z_2\).
    \task Stellen Sie \(z_2\) in der exponentiellen Form dar.
\end{tasks}
\subsection*{Lösung}

\newpage
\begin{center}
    {\LARGE \textbf{Credits}}
\end{center}
\vspace{1cm}
\noindent\textbf{Main Editor: } Milix
\vspace{0.5cm}
\begin{center}
    \href{https://github.com/Milix244}{GITHUB PROFILE}  \\
    \vspace{0.3cm}
    \href{https://github.com/Milix244/Reserve-BAC-2017}{GITHUB REPO}  \\
    \vspace{0.3cm}
    \href{https://discord.gg/EpQ3eukxJ4}{DISCORD SERVER}
\end{center}
\vspace{1cm}
\textit{\copyright\ 2026 -- Alle Rechte vorbehalten trust.}


\end{document}